\section{Odtworzenie wyników}
\label{app:reproduction}

Do odtworzenia wyników potrzebne są kompilator języka Rust w wersji 1.88.0, wskazanej w pliku \texttt{rust-toolchain.toml}, oraz Python w wersji co najmniej 3.13 z menedżerem pakietów \texttt{uv}, który instaluje zależności w wersjach zapisanych w pliku \texttt{uv.lock}. Ścieżki w plikach konfiguracyjnych są względne, dlatego polecenia uruchamia się z katalogu głównego repozytorium. Polecenia przedstawiono na Listingu~\ref{lst:reproduction}.

\begin{code}
\begin{Verbatim}[fontsize=\small]
# zbudowanie programu i instalacja zależności Pythona
cargo build --release --manifest-path src/compute/Cargo.toml
uv sync --project src/analysis

# eksperyment pokazowy: uczenie i wykresy
cargo run --release --manifest-path src/compute/Cargo.toml -- \
    train data/experiments/classification/showcase-three-gauss
uv run --project src/analysis src/analysis/main.py \
    plot data/experiments/classification/showcase-three-gauss

# jedno badanie: eksperymenty, uczenie, zestawienie wyników
cargo run --release --manifest-path src/compute/Cargo.toml -- \
    generate data/experiments/classification/simple-activations
cargo run --release --manifest-path src/compute/Cargo.toml -- \
    train data/experiments/classification/simple-activations
uv run --project src/analysis src/analysis/main.py \
    aggregate data/experiments/classification/simple-activations --plot-runs

# wszystkie badania części pierwszej (PowerShell 7)
pwsh -File scripts/run-studies.ps1
\end{Verbatim}
    \caption{Polecenia odtwarzające wyniki, uruchamiane z katalogu głównego repozytorium.}
    \label{lst:reproduction}
\end{code}

Konfiguracje i wyniki znajdują się w następujących miejscach:
\begin{itemize}
    \item konfiguracje eksperymentów pokazowych:\newline \texttt{data/experiments/<zadanie>/showcase-<zbiór>/config.toml},
    \item konfiguracje badań:\newline \texttt{data/experiments/<zadanie>/<zbiór>-<czynnik>/study.toml},
    \item wykresy eksperymentów pokazowych: katalog \texttt{plots/} eksperymentu,
    \item zestawienia badań: katalog \texttt{results/} badania.
\end{itemize}
Katalog \texttt{<zadanie>} to \texttt{classification} albo \texttt{regression}. Eksperymenty pokazowe używają ziarna $1$, a badania ziaren od $1$ do $10$.
